% SPDX-License-Identifier: CC-BY-SA-4.0
% Author: Matthieu Perrin

\begingroup

\begin{exercice}[Langages engendrés par des grammaires rationnelles]\label{exo:parsing/rational/rational_to_other}

  Traduire chaque grammaire en une grammaire rationnelle à droite, en un automate à états finis, et en une expression rationnelle.

  \begin{question}
  \item $G_1 \eqdef \left\langle\{cha, bada\}, \{S, A\}, S, \left\{
    \begin{array}{lll}
      S   &\rightarrow& cha bada A \\
      A   &\rightarrow& bada A \mid bada
    \end{array}\right\}\right\rangle$
  \item $G_2 \eqdef \left\langle\{a, b\}, \{S\}, S, \left\{
    \begin{array}{lll}
      S   &\rightarrow& \varepsilon \;|\; Sa \;|\; b 
    \end{array}\right\}\right\rangle$
  \item $G_3 \eqdef \left\langle\{a, b, c\}, \{S, A, B\}, S, \left\{
    \begin{array}{lll}
      S   &\rightarrow& \varepsilon \;|\; A ,\\
      A   &\rightarrow& aA \;|\; B \;|\; bS ,\\
      B   &\rightarrow& \varepsilon \;|\; cA \;|\; cB
    \end{array}\right\}\right\rangle$
  \end{question}

\end{exercice}

\begin{correction}

  \begin{question}

  \item Pour $G_1$ :
    \begin{itemize}
    \item Grammaire rationnelle à droite :
      \[
        \left\langle\{cha,bada\},\{S,A,F\},S,
        \left\{\begin{array}{rcl}
          S &\rightarrow& cha\,F\\
          F &\rightarrow& bada\,A\\
          A &\rightarrow& bada\,A \mid bada
        \end{array}\right\}\right\rangle
        \]
      \item Automate fini :
          \begin{tikzpicture}[automaton, baseline=(S.base), x=20mm]
            \state[initial]   (S) at (0,0) {$S$};
            \state            (F) at (1,0) {$F$};
            \state            (A) at (2,0) {$A$};
            \state[accepting] (T) at (3,0) {$T$};

            \path (S) edge node {$cha$} (F);
            \path (F) edge node {$bada$} (A);
            \path (A) edge[loop above] node {$bada$} (A);
            \path (A) edge node {$bada$} (T);
          \end{tikzpicture}
      \item Expression rationnelle : $cha\;bada^+$.
    \end{itemize}

  \item Pour $G_2$ :
    \begin{itemize}
    \item Grammaire rationnelle à droite :
      \[
        \left\langle\{a,b\},\{S,F\},S,
        \left\{\begin{array}{rcl}
          S &\rightarrow& bT \mid \varepsilon \\
          T &\rightarrow& aT \mid \varepsilon
        \end{array}\right\}\right\rangle
      \]
    \item Automate fini :
        \begin{tikzpicture}[automaton, baseline=(S.base)]
          \state[initial, accepting] (S) at (0,0) {$S$};
          \state[accepting]          (T) at (1,0) {$T$};

          \path (S) edge node {$b$} (T);
          \path (T) edge[loop above] node {$a$} (T);
        \end{tikzpicture}
    \item Expression rationnelle : $ba^\star\mid \varepsilon$.
    \end{itemize}

  \item Pour $G_3$ :
    \begin{itemize}
    \item Grammaire rationnelle à droite :
      \[
        \left\langle\{a,b,c\},\{S,A,B\},S,
        \left\{\begin{array}{rcl}
          S &\rightarrow& aA \mid bA \mid cA \mid cB \mid a \mid b \mid c \mid \varepsilon\\
          A &\rightarrow& aA \mid bA \mid cA \mid cB \mid a \mid b \mid c\\
          B &\rightarrow& cA \mid cB \mid c
        \end{array}\right\}\right\rangle
      \]
    \item Automate fini :
      \begin{tikzpicture}[automaton, baseline=(S.base), size=20mm]
        \state[initial, accepting] (S) at (0,1) {$S$};
        \state                     (A) at (2,2) {$A$};
        \state                     (B) at (2,0) {$B$};
        \state[accepting]          (F) at (1,1) {$F$};

        \path (S) edge             node[sloped] {$a, b, c$} (A);
        \path (S) edge             node[sloped] {$c$}       (B);
        \path (S) edge             node[sloped] {$a, b, c$} (F);
        \path (A) edge[loop above] node[sloped] {$a, b, c$} (A);
        \path (A) edge[bend left]  node[sloped] {$c$} (B);
        \path (A) edge             node[sloped] {$a, b, c$} (F);
        \path (B) edge[loop below] node[sloped] {$c$} (B);
        \path (B) edge[bend left]  node[sloped] {$c$} (A);
        \path (B) edge             node[sloped] {$c$} (F);
      \end{tikzpicture}
    \item Expression rationnelle : $(a\mid b \mid c)^\star$.
    \end{itemize}

  \end{question}

\end{correction}

\endgroup
\endinput
